\subsection{道路的提升}

命题 3. --- 设 B 是一个拓扑空间，\(p\colon E \to B\) 是一个覆叠。记 \(\widetilde{p}\colon \mathcal{C}_{c}(\mathbf{I};E) \to \mathcal{C}_{c}(\mathbf{I};B)\) 为映射 \(c \mapsto p \circ c\)。记 \(o_{\mathrm{E}}\colon \mathcal{C}_{c}(\mathbf{I};E) \to E\)（分别记 \(o_{\mathrm{B}}\colon \mathcal{C}_{c}(\mathbf{I};B) \to B\)）为由 \(c \mapsto c(0)\) 定义的映射。则图

\[
\begin{array}{c} \mathcal {C} _ {c} (\mathbf {I}; \mathrm{E}) \xrightarrow {o _ {\mathrm{E}}} \mathrm{E} \\ \Big \downarrow \widetilde {p} \\ \mathcal {C} _ {c} (\mathbf {I}; \mathrm{B}) \xrightarrow {o _ {\mathrm{B}}} \mathrm{B} \end{array}
\]

是一个笛卡尔方块。

拓扑空间 I 局部连通、局部紧且单连通。因此，将 I, p.~131 的命题 9 应用于 T = I、t = 0，即得本命题。

推论 1. --- 映射 \(\widetilde{p}\colon\mathcal{C}_{c}(\mathbf{I};\mathrm{E})\to\mathcal{C}_{c}(\mathbf{I};\mathrm{B})\) 是一个覆叠。

这由 I, p.~71 的命题 1 的推论 2 得到。

推论 2（道路的提升）. --- 设 \(p\) : E → B 是一个覆叠，x 是 E 中的一点，a = p(x)。映射 c ↦ p ∘ c 是从 E 中起点为 x 的道路空间 \(\Lambda_x\) (E) 到 B 中起点为 a 的道路空间 \(\Lambda_a\) (B) 的同胚。

用命题 3 的记号，有 \(\Lambda_{a}(\mathrm{B}) = o_{\mathrm{B}}^{-1}(a)\) 和 \(\Lambda_x(\mathrm{E}) = o_{\mathrm{E}}^{-1}(x)\)。推论于是由 I, p.~10 的命题 4 的推论得到。

如果连续映射 \(p \colon \mathrm{E} \to \mathrm{B}\) 满足下述道路提升性质，我们就这样称呼它：对每条道路 \(c \colon \mathbf{I} \to \mathbf{B}\) 以及 E 中提升 \(x\) 的每个点 \(c(0)\)，存在一条 E 中起点为 \(c'\) 且提升 c 的道路 \(x\)。若 \(c\)、\(p\) 是平展且分离的，这样的道路便是唯一的（I, p.~34, 命题 11 的推论 1）。设 E 是一个覆叠，\(p\) 是其投影；映射 \(p\) 是平展、分离的，并且具有道路提升性质（推论 2）。部分逆命题见 III, p.~315 的推论。

命题 4. --- 设 \(p \colon \mathrm{E} \to \mathrm{B}\) 是满足道路提升性质的平展分离映射。设 \(a\) 和 \(b\) 是 \(\mathrm{B}\) 中的点，\(x\) 是 \(\mathrm{E}\) 中满足 \(p(x) = a\)、\(c_0\)、\(c_1\) 的点。两条在 \(a\) 中连接 \(b\) 与 \(\mathrm{B}\) 的严格同伦道路，其分别以 x 为起点在 \(x\)、\(\mathrm{E}\)、\(c_0\)、\(c_1\) 中的提升具有相同的终点，并且严格同伦。

设 \(\sigma\colon\mathbf{I}\times\mathbf{I}\to\mathbf{B}\) 是连接 \(c_{0}\) 与 \(c_{1}\) 的严格同伦。对 \(s\in\mathbf{I}\)，令 \(c_{s}^{\prime}\) 为 E 中以 x 为起点、提升道路 \(x\)、\(t\mapsto\sigma(t,s)\) 的唯一道路。对 \((t,s)\in\mathbf{I}\times\mathbf{I}\)，置 \(\sigma^{\prime}(t,s)=c_{s}^{\prime}(t)\)；有 \(p\circ\sigma^{\prime}=\sigma\)。映射 \(\sigma^{\prime}\) 在 \(\{0\}\times\mathbf{I}\) 上为常值；按构造，它在每个 \(\mathbf{I}\times\{s\}\) 上连续。因此由 I, p.~37, 定理 1 的推论 1，它是连续的。特别地，从 \(s\in\mathbf{I}\)、\(\mathbf{I}\) 到 \(\mathrm{E}_{b}\) 的映射 \(s\mapsto\sigma^{\prime}(1,s)\)、\(b\) 是提升像为 b 的常值道路的连续映射；因而它是常值的。这证明映射 \(\sigma^{\prime}\) 是连接 \(c_{0}^{\prime}\) 与 \(c_{1}^{\prime}\) 的严格同伦。

推论 1. --- 设 B 是一个拓扑空间，E 是 B 的一个覆叠；记 p 为其投影。对 E 中的每一对 \((x,y)\)，映射 \(\varpi_{x,y}(p)\colon\varpi_{x,y}(\mathrm{E})\to\varpi_{p(x),p(y)}(\mathrm{B})\) 是单射。特别地，对 E 中的每个点 x，同态 \(\pi_{1}(p,x)\colon\pi_{1}(\mathrm{E},x)\to\pi_{1}(\mathrm{B},p(x))\) 是单射。

推论 2. --- 设 B 是一个拓扑空间，E 是 B 的一个覆叠；记 p 为其投影。令 x 为 E 中的一点，置 \(b = p(x)\)。对于 B 中 a 处的圈 c，其在 \(\pi_{1}(B, b)\) 中的类属于同态 \(\pi_{1}(p, x)\) 的像子群，当且仅当以 x 为起点提升 c 的道路 c' 是 E 中 x 处的圈。

这个条件显然是充分的。反之，设 \(c_{1}^{\prime}\) 是 E 中 x 处的一个圈。置 \(c_{1} = p \circ c_{1}^{\prime}\)，并假设圈 c 与 \(c_{1}\) 严格同伦。由命题 4，道路 \(c^{\prime}\) 与道路 \(c_{1}^{\prime}\) 具有相同的终点。因此它是 x 处的圈。

推论 3. --- 设 B 是一个拓扑空间，\((\mathrm{E}_{1}, p_{1})\) 和 \((\mathrm{E}_{2}, p_{2})\) 是 B 的覆叠。记 \((\mathrm{E}, p)\) 为它们的 B-积空间，\(x = (x_{1}, x_{2})\) 是 E 中的一点。同态 \(\pi_{1}(p, x)\) 的像等于同态 \(\pi_{1}(p_{1}, x_{1})\) 和 \(\pi_{1}(p_{2}, x_{2})\) 的像的交。

设 \(c\) 是 B 中以 \(p_{1}(x_{1}) = p_{2}(x_{2})\) 为基点的圈，并且对 \(i \in \{1, 2\}\)，令 \(c_{i}^{\prime}\) 为 \(x_i\) 中以 \(E_i\) 为起点、提升 c 的道路。B-空间 \(c\)、\(E = E_{1} \times_{B} E_{2}\) 是 B 的一个覆叠（I, p.~72, 命题 2 的推论 3），而道路 \(c^{\prime}: t \mapsto (c_{1}^{\prime}(t), c_{2}^{\prime}(t))\) 是 E 中以 x 为起点提升 c 的唯一道路。\(x\)、\(c\)、\(c^{\prime}\) 是圈的充要条件是 \(c_{1}^{\prime}\) 和 \(c_{2}^{\prime}\) 都是圈。因此推论 3 由推论 2 得到。

